\begin{abcliste}
\abc Die molare Masse von Wasser ist (FoTaBe S.~228)
\begin{align}
 M &= \MWdF = 2\cdot\MHaO + \MOaO = \MWdg
\end{align}
Mit $T = \TDF = \TD$ und $\bar v = \sqrt{8kT/(\pi m)} = \sqrt{8RT/(\pi M)}$ (FoTaBe S.~171) wird
\begin{align}
 \bar v &= \vDF = \sqrt{\frac{8\cdot\ncR\cdot\TD}{\pi\cdot\MWd}} = \vD \approx \result{\vDP{0}{2}}
\end{align}

\abc Der Dampf verhält sich näherungsweise wie ein ideales Gas, $pV = \frac{m}{M}RT$:
\begin{align}
 p &= \pDF = \frac{\mWa\cdot\ncR\cdot\TD}{\MWd\cdot\VD} = \pD \approx \result{\pDbP{0}{2}}
\end{align}
Kontrolle: Die Dichte des Dampfes ist $\rho = \rDF = \rD$. Das ist weniger als die Dichte des gesättigten Wasserdampfs bei \thDO, $\rho_\text{s} = \rSa$, und der Druck ist kleiner als der Sättigungsdampfdruck $p_\text{s} = \pSbP{0}{3}$ (FoTaBe S.~192). Das Wasser kann also tatsächlich vollständig verdampfen.
\end{abcliste}
